\chapter{解析函数}
	
	\section{复变函数的极限}
	设函数 \(w=f(z)=u(x,y)+iv(x,y)\) 在 \(z_0=x_0+iy_0\) 的去心邻域内有定义，\(A=a+ib\) 为复常数。若对任意 \(\varepsilon>0\)，存在 \(\delta>0\)，使得当
	\[
	0<|z-z_0|<\delta
	\]
	时，有
	\[
	|f(z)-A|<\varepsilon,
	\]
	则称 \(A\) 为 \(f(z)\) 当 \(z\to z_0\) 时的极限，记作
	\[
	\lim_{z\to z_0}f(z)=A.
	\]
	
	\subsection{极限存在的充要条件}
	复变函数极限存在等价于其实部和虚部的极限同时存在：
	\[
	\lim_{z\to z_0}f(z)=A\iff
	\lim_{(x,y)\to(x_0,y_0)}u(x,y)=a,\quad
	\lim_{(x,y)\to(x_0,y_0)}v(x,y)=b.
	\]
	
	\subsection{极限运算法则}
	若 \(\lim\limits_{z\to z_0}f(z)=A\)，\(\lim\limits_{z\to z_0}g(z)=B\)，则
	\begin{enumerate}
		\item \(\lim\limits_{z\to z_0}[f(z)\pm g(z)]=A\pm B\);
		\item \(\lim\limits_{z\to z_0}[f(z)\cdot g(z)]=AB\);
		\item \(\lim\limits_{z\to z_0}\dfrac{f(z)}{g(z)}=\dfrac{A}{B}\ (B\neq 0)\).
	\end{enumerate}
	
	\section{复变函数的连续性}
	\begin{tcolorbox}[title=定义：复变函数的连续性, colframe=blue!50!black, colback=blue!5!white, sharp corners]
		若 \(\lim\limits_{z\to z_0}f(z)=f(z_0)\)，则称 \(f(z)\) 在 \(z_0\) 处\textbf{连续}。
	\end{tcolorbox}
	
	\subsection{连续的充要条件}
	\[
	f(z)\ \text{在}\ z_0\ \text{连续}
	\iff
	u(x,y),\ v(x,y)\ \text{在}\ (x_0,y_0)\ \text{处均连续}.
	\]
	
	\subsection{连续函数的运算性质}
	\begin{enumerate}
		\item 连续函数的和、差、积、商（分母非零）仍连续；
		\item 连续函数的复合函数仍连续；
		\item 有界闭域上的连续函数必有界、可取到最值。
	\end{enumerate}
	
	\section{复变函数的导数}
	\begin{tcolorbox}[title=定义：复变函数的导数, colframe=blue!50!black, colback=blue!5!white, sharp corners]
		设函数 \(w=f(z)\) 在 \(z_0\) 的邻域内有定义。若极限
		\[
		\lim_{\Delta z\to 0}\frac{\Delta w}{\Delta z}
		=\lim_{\Delta z\to 0}\frac{f(z_0+\Delta z)-f(z_0)}{\Delta z}
		\]
		存在且有限，则称 \(f(z)\) 在 \(z_0\) 处\textbf{可导}，该极限称为 \(f(z)\) 在 \(z_0\) 处的导数，记作
		\[
		f'(z_0)=\left.\frac{\mathrm{d}w}{\mathrm{d}z}\right|_{z=z_0}.
		\]
	\end{tcolorbox}
	
	\subsection{可导的必要条件：连续}
	若 \(f(z)\) 在 \(z_0\) 可导，则 \(f(z)\) 必在 \(z_0\) 连续。
	\[
	\lim_{\Delta z\to0}\big[f(z_0+\Delta z)-f(z_0)\big]
	=\lim_{\Delta z\to0}\frac{f(z_0+\Delta z)-f(z_0)}{\Delta z}\cdot\Delta z
	=f'(z_0)\cdot 0=0.
	\]
	
	\section{柯西—黎曼方程（C-R 方程）}
	设 \(f(z)=u(x,y)+iv(x,y)\)，\(z=x+iy\)。
	
	\begin{tcolorbox}[title=定理：柯西—黎曼方程 (C-R 方程), colframe=green!45!black, colback=green!5!white]
		\(f(z)\) 在 \(z\) 处可导的必要条件是 \(u,v\) 可微且满足
		\[
		\frac{\partial u}{\partial x}=\frac{\partial v}{\partial y},\qquad
		\frac{\partial u}{\partial y}=-\frac{\partial v}{\partial x}.
		\]
		若 \(u,v\) 一阶偏导连续且满足 C-R 方程，则 \(f(z)\) 可导。
	\end{tcolorbox}
	
	\subsection{推导思路}
	令 \(\Delta z=\Delta x+i\Delta y\)，分别沿实轴、虚轴取极限，令两者相等即得 C-R 方程。
	
	\subsection{导数计算公式}
	\[
	f'(z)=\frac{\partial u}{\partial x}+i\frac{\partial v}{\partial x}
	=\frac{\partial v}{\partial y}-i\frac{\partial u}{\partial y}.
	\]
	
	\section{解析函数}
	\begin{tcolorbox}[title=定义：解析函数, colframe=blue!50!black, colback=blue!5!white, sharp corners]
		若 \(f(z)\) 在 \(z_0\) 及其某一邻域内\textbf{处处可导}，则称 \(f(z)\) 在 \(z_0\) 处\textbf{解析}。
		
		若 \(f(z)\) 在区域 \(D\) 内每一点都解析，则称 \(f(z)\) 是 \(D\) 内的\textbf{解析函数}。
	\end{tcolorbox}
	
	\subsection{可导与解析的区别}
	\begin{tcolorbox}[colframe=orange!75!black, colback=orange!5!white, leftrule=4pt, rightrule=0pt, toprule=0pt, bottomrule=0pt, sharp corners]
		\begin{itemize}
			\item \textbf{可导}：只要求在一点存在导数；
			\item \textbf{解析}：要求在一点及其邻域内都可导，条件更强。
		\end{itemize}
	\end{tcolorbox}
	
	\subsection{解析的充要条件}
	\(f(z)=u+iv\) 在区域 \(D\) 内解析 \(\iff\)
	\begin{enumerate}
		\item \(u(x,y),\ v(x,y)\) 在 \(D\) 内可微；
		\item \(u,v\) 在 \(D\) 内处处满足 C-R 方程。
	\end{enumerate}
	
	\section{调和函数}
	\begin{tcolorbox}[title=定义：调和函数, colframe=blue!50!black, colback=blue!5!white, sharp corners]
		若二元实函数 \(u(x,y)\) 在区域 \(D\) 内具有二阶连续偏导数，且满足拉普拉斯方程
		\[
		\Delta u=\frac{\partial^2 u}{\partial x^2}+\frac{\partial^2 u}{\partial y^2}=0,
		\]
		则称 \(u\) 为 \(D\) 内的\textbf{调和函数}。
	\end{tcolorbox}
	
	\subsection{解析函数与调和函数的关系}
	\begin{tcolorbox}[title=定理：解析函数与调和函数, colframe=green!45!black, colback=green!5!white]
		若 \(f(z)=u+iv\) 在区域 \(D\) 内解析，则 \(u,v\) 都是 \(D\) 内的调和函数。
	\end{tcolorbox}
	
	\textbf{推导：}
	由解析性知 C-R 方程成立：
	\[
	u_x=v_y,\quad u_y=-v_x.
	\]
	分别对两式求导：
	\[
	u_{xx}=v_{yx},\quad u_{yy}=-v_{xy}.
	\]
	由二阶偏导连续，\(v_{yx}=v_{xy}\)，故
	\[
	u_{xx}+u_{yy}=0.
	\]
	同理可证 \(v_{xx}+v_{yy}=0\)。
	
	\subsection{共轭调和函数}
	\begin{tcolorbox}[title=定义：共轭调和函数, colframe=blue!50!black, colback=blue!5!white, sharp corners]
		设 \(u,v\) 均为调和函数，且满足 C-R 方程，则称 \(v\) 为 \(u\) 的\textbf{共轭调和函数}。
	\end{tcolorbox}
	
	若 \(f(z)=u+iv\) 解析，则 \(v\) 是 \(u\) 的共轭调和函数。已知 \(u\) 可通过 C-R 方程求出 \(v\)，从而构造解析函数。
	
